% Avance de Trabajo de Investigación 2 (ATI2): Metodología de la RSL
% Plantilla: Springer LNCS (clase llncs, disponible por defecto en Overleaf)
\documentclass[runningheads]{llncs}
\usepackage[T1]{fontenc}
\usepackage[utf8]{inputenc}
\usepackage{lmodern}
\usepackage[spanish,es-noshorthands,es-tabla]{babel}
\usepackage{graphicx}
\usepackage{booktabs}
\usepackage{longtable}
\usepackage{array}
\usepackage{tabularx}
\usepackage{tikz}
\usetikzlibrary{positioning,arrows.meta,calc}
\usepackage{url}
\usepackage[hidelinks]{hyperref}
\renewcommand\UrlFont{\color{blue}\rmfamily}
\newcolumntype{L}[1]{>{\raggedright\arraybackslash}p{#1}}

\begin{document}

\title{Integración de IA Generativa Multimodal en la Interacción Humano-Robot (HRI) para Robótica Social (2020--2026)}
\titlerunning{IA Generativa Multimodal en HRI para Robótica Social}

\author{Alejandro Valentino Seminario Medina\orcidID{0009-0006-5378-9115} \and
Alvaro Jordy Ordaya Zapata\orcidID{0009-0005-0751-5474}}
\authorrunning{A. V. Seminario Medina y A. J. Ordaya Zapata}
\institute{Universidad Tecnológica del Perú, Piura, Perú\\
\email{U22247454@utp.edu.pe, U20303695@utp.edu.pe}}

\maketitle

\begin{abstract}
La incorporación de modelos generativos multimodales en la robótica social transformó la interacción humano-robot entre 2020 y 2026, pero la evidencia documenta limitaciones técnicas (latencia, errores de reconocimiento de voz y desincronización entre habla y gesto), cognitivas (fragilidad de la Teoría de la Mente y del razonamiento espacial) y éticas (sobre-antropomorfización y vulnerabilidad a la manipulación conversacional). Este avance presenta la metodología de una revisión sistemática de literatura (RSL) orientada a sistematizar las arquitecturas de fusión sensorial y de traducción de lenguaje a acción que abordan estas limitaciones. Se formuló una pregunta PICO, se definieron palabras clave controladas y una ecuación de búsqueda booleana que se aplicó de forma equivalente en Semantic Scholar y Crossref, y se seleccionaron los estudios conforme a la lógica de la declaración PRISMA 2020. De 942 registros identificados se eliminaron 64 duplicados, se cribaron 592 registros por título y resumen, se buscaron 95 informes a texto completo, se leyeron los 37 recuperables y se incluyeron 22 estudios primarios publicados entre 2023 y 2026.

\keywords{Robótica social \and Interacción humano-robot \and IA generativa multimodal \and Robots socialmente asistidos \and Revisión sistemática \and PRISMA.}
\end{abstract}

%=====================================================================
\section{Introducción}
%=====================================================================
En el campo de la interacción humano-robot (HRI) aplicada a la robótica social, la integración de la inteligencia artificial generativa multimodal ha cobrado relevancia entre 2020 y 2026. La aparición de los modelos de lenguaje de gran tamaño (LLM) y de arquitecturas capaces de procesar texto, voz, imágenes y video de forma conjunta desplazó el diálogo basado en gramáticas cerradas y árboles de decisión por respuestas contextuales no anticipadas por el diseñador. Zhang y Soh (2023) señalan que un LLM puede operar como modelo humano \textit{zero-shot}; Kim et al. (2024) muestran que estos modelos reconfiguran las expectativas de diseñadores y usuarios; y Lai et al. (2025) evidencian, con la arquitectura NVP-HRI, la viabilidad de fusionar voz y postura en plataformas físicas.

No obstante, la literatura presenta limitaciones. En el plano técnico, la memoria externa mediante generación aumentada por recuperación introduce retrasos de entre seis y quince segundos (Mauliana et al., 2025), los errores del reconocimiento automático de voz se concentran en hablantes no nativos y el modelo los encubre con respuestas plausibles (Verhelst y Belpaeme, 2025), y la voz expresiva y la gesticulación se generan en flujos independientes (Barakat et al., 2024). En el plano cognitivo, la Teoría de la Mente atribuida a los LLM se degrada ante perturbaciones triviales (Verma et al., 2024) y la falta de razonamiento espacial induce planes motores inseguros (Lai et al., 2025). En el plano ético, la expresividad del robot propicia la sobre-antropomorfización (Abercrombie et al., 2023) y los usuarios recurren al \textit{gaslighting} o al \textit{roleplay} para vulnerar sus resguardos (Abbo et al., 2025). Las revisiones previas abordan estos aspectos por separado (Barakat et al., 2024; Maure y Bruno, 2025) y no cubren la evidencia de 2024--2026.

El objetivo de la RSL es sistematizar el estado del arte sobre la integración de IA generativa multimodal en la HRI para robótica social entre 2020 y 2026, con énfasis en las arquitecturas propuestas para mitigar la latencia, sincronizar la comunicación no verbal y garantizar la gobernanza ética de la interacción. El presente avance describe la metodología seguida para identificar y seleccionar los estudios.

%=====================================================================
\section{Metodología}
%=====================================================================
\subsection{Tipo de estudio y protocolo}
Se realizó una revisión sistemática de literatura siguiendo las directrices de Kitchenham y Charters (2007) para ingeniería de software y la lógica de selección y reporte de la declaración PRISMA 2020 (Page et al., 2021). El proceso comprendió cinco fases: (1) formulación de la pregunta mediante la estrategia PICO; (2) definición de palabras clave y construcción de la ecuación de búsqueda; (3) definición de los criterios de elegibilidad; (4) búsqueda en bases de datos científicas y exportación de registros a Excel; y (5) selección sistemática de estudios documentada en un diagrama de flujo PRISMA. La búsqueda, la eliminación de duplicados y la aplicación de filtros se ejecutaron mediante scripts reproducibles que acompañan a la base de datos en Excel.

\subsection{Pregunta de investigación (estrategia PICO)}
La estrategia PICO permite descomponer la pregunta de investigación en componentes que se traducen directamente en términos de búsqueda (Santos et al., 2007). La Tabla~\ref{tab:pico} presenta los componentes y sus palabras clave. La pregunta formulada es:

\begin{quote}
\textit{En robots sociales, socialmente asistivos y de compañía (P), ¿qué arquitecturas de integración de IA generativa multimodal ---fusión sensorial y traducción de lenguaje a acción con LLM o VLM--- (I), frente a los enfoques de diálogo guionado o unimodal (C), se han propuesto y qué resultados reportan sobre la latencia de respuesta, la sincronización de la comunicación no verbal, el desempeño cognitivo y la gobernanza ética de la interacción (O)?}
\end{quote}

\begin{table}[htbp]
\centering
\caption{Componentes de la pregunta PICO y palabras clave asociadas.}
\label{tab:pico}
\footnotesize
\begin{tabular}{L{1.2cm}L{3.5cm}L{6.6cm}}
\toprule
\textbf{Comp.} & \textbf{Descripción} & \textbf{Palabras clave (inglés)} \\
\midrule
P & Robots sociales, socialmente asistivos y de compañía en HRI & ``social robot'', ``social robotics'', ``socially assistive robot'', ``companion robot'', ``human-robot interaction'' \\
I & IA generativa multimodal: LLM, VLM y modelos fundacionales integrados al robot & ``large language model'', ``generative AI'', ``generative artificial intelligence'', ``vision-language model'', ``foundation model'', ``GPT'', ``ChatGPT'', ``multimodal LLM'' \\
C & Diálogo guionado, basado en reglas o unimodal & No se incluye en la ecuación (componente no esencial que restringe en exceso los resultados) \\
O & Latencia, sincronización verbal--no verbal, ToM y razonamiento espacial, confianza, antropomorfización, manipulación y seguridad & Se evalúa en el cribado (criterio CI4), no en la ecuación \\
\bottomrule
\end{tabular}
\end{table}

\subsection{Palabras clave}
Las palabras clave se obtuvieron de la ficha del tema de investigación y de los términos empleados en los títulos, resúmenes y descriptores de los estudios de referencia (Kim et al., 2024; Lai et al., 2025; Maure y Bruno, 2025). La Tabla~\ref{tab:kw} muestra los términos principales en español e inglés y sus sinónimos.

\begin{table}[htbp]
\centering
\caption{Palabras clave especializadas y sinónimos.}
\label{tab:kw}
\footnotesize
\begin{tabular}{L{3.3cm}L{3.3cm}L{4.8cm}}
\toprule
\textbf{Español} & \textbf{Inglés} & \textbf{Sinónimos / variantes} \\
\midrule
Robótica social & Social robotics & social robot(s) \\
Robots socialmente asistivos & Socially assistive robots & companion robot(s) \\
Interacción humano-robot & Human-robot interaction & HRI \\
IA generativa multimodal & Multimodal generative AI & generative AI, multimodal LLM, vision-language model \\
Modelos de lenguaje de gran tamaño & Large language models & LLM, GPT, ChatGPT, foundation model \\
\bottomrule
\end{tabular}
\end{table}

\subsection{Ecuación de búsqueda}
Los sinónimos de cada componente se agruparon con el operador \texttt{OR}, y los componentes P e I se unieron con \texttt{AND}. Los componentes C y O se excluyeron de la ecuación porque su inclusión reducía los resultados a conjuntos truncos. Se aplicaron en el cribado mediante los criterios de elegibilidad. La ecuación genérica es:

\begin{quote}\small\ttfamily\raggedright
("social robot*" OR "social robotics" OR "socially assistive robot*" OR "companion robot*" OR "human-robot interaction") AND ("large language model*" OR "generative AI" OR "generative artificial intelligence" OR "vision-language model*" OR "foundation model*" OR GPT OR ChatGPT OR "multimodal LLM")
\end{quote}

Esta ecuación canónica se aplica sin cambios en todas las bases de datos; solo varía el campo de búsqueda y la forma de aplicar los filtros (Tabla~\ref{tab:eq}):

\begin{itemize}
\item \textbf{Scopus.} La ecuación se pega en la caja \textit{Search documents} con el campo \textit{Article title, Abstract, Keywords}. Se aplican los filtros \textit{Year range} 2020--2026, \textit{Document type} Article y Conference Paper, y \textit{Language} English y Spanish. Su forma en la búsqueda avanzada es:\\
{\ttfamily\raggedright TITLE-ABS-KEY(<ecuación>) AND PUBYEAR > 2019 AND PUBYEAR < 2027 AND (LIMIT-TO(DOCTYPE,"ar") OR LIMIT-TO(DOCTYPE,"cp")) AND (LIMIT-TO(LANGUAGE,"English") OR LIMIT-TO(LANGUAGE,"Spanish"))\par}
\item \textbf{Web of Science Core Collection.} En la \textit{Fielded Search} se elige el campo \textit{Topic} y se pega la ecuación tal cual, \emph{sin} el prefijo \texttt{TS=}, porque el campo ya lo aplica. Se fija \textit{Publication Date} del 2020-01-01 al 2026-12-31. Su equivalente en el \textit{Query Builder} es \texttt{TS=(<P>) AND TS=(<I>)}.
\item \textbf{Semantic Scholar.} El operador \texttt{+} equivale a \texttt{AND} y \texttt{|} a \texttt{OR}. Las formas plurales se escriben explícitamente porque no admite truncamiento con \texttt{*}. La búsqueda se hace en título y resumen, con años 2020--2026.
\item \textbf{Crossref.} Su API no admite operadores booleanos, por lo que se ejecutaron las 30 combinaciones P$\times$I. Después se aplicó la misma ecuación sobre el título y el resumen de cada registro, de modo que el resultado es lógicamente equivalente.
\end{itemize}

\begin{table}[htbp]
\centering
\caption{Aplicación de la ecuación canónica por base de datos.}
\label{tab:eq}
\scriptsize
\begin{tabular}{L{1.9cm}L{2.6cm}L{4.6cm}L{1.9cm}}
\toprule
\textbf{Base} & \textbf{Campo} & \textbf{Filtros} & \textbf{Registros} \\
\midrule
Semantic Scholar & Título y resumen (API \textit{bulk search}) & 2020--2026; ejecución 3/10/2026 & 873 \\
\addlinespace
Crossref & Título y resumen (API REST, filtro booleano local) & 2020--2026; ejecución 3/10/2026 & 69 \\
\addlinespace
Scopus & Article title, Abstract, Keywords & 2020--2026; Article y Conference Paper; English y Spanish & Por ejecutar \\
\addlinespace
WoS Core Collection & Topic (Fielded Search, sin \texttt{TS=}) & Publication Date 2020-01-01 a 2026-12-31 & Por ejecutar \\
\midrule
\multicolumn{3}{r}{\textbf{Total identificado}} & \textbf{942} \\
\bottomrule
\end{tabular}
\end{table}

\subsection{Criterios de inclusión y exclusión}
Los criterios de elegibilidad se derivaron de los componentes de la pregunta PICO (Tabla~\ref{tab:criterios}) y se aplicaron tanto en el cribado por título y resumen como en la evaluación de elegibilidad. El periodo 2020--2026 no es arbitrario. Comienza con la publicación de los primeros LLM generativos de propósito general, que habilitaron la conversación abierta en robots sociales, y cubre la evidencia posterior a la ventana de las revisiones previas (Barakat et al., 2024; Maure y Bruno, 2025).

\begin{table}[htbp]
\centering
\caption{Criterios de inclusión (CI) y de exclusión (CE).}
\label{tab:criterios}
\footnotesize
\begin{tabular}{L{0.9cm}L{10.4cm}}
\toprule
\textbf{Cód.} & \textbf{Criterio} \\
\midrule
CI1 & El estudio aborda robots sociales, socialmente asistivos o de compañía, o la HRI con un robot físico (P). \\
CI2 & El robot integra IA generativa: LLM, VLM, modelos fundacionales o GPT (I). \\
CI3 & La interacción es multimodal: combina al menos dos canales entre voz, texto, imagen, video, gesto, mirada, expresión facial o señales fisiológicas (I). \\
CI4 & Reporta una evaluación empírica (estudio con usuarios, experimento, piloto o evaluación técnica) con resultados sobre latencia, comportamiento no verbal, desempeño cognitivo, aceptación, confianza o aspectos éticos (O). \\
CI5 & Artículo de revista o de conferencia revisado por pares, publicado entre 2020 y 2026, en inglés o español. \\
\midrule
CE1 & Preprints sin revisión por pares, tesis, libros, capítulos, editoriales y otros documentos no indexados. \\
CE2 & Robots sin función social o asistencial (conducción autónoma, manipulación industrial, logística o tareas domésticas sin interacción social). \\
CE3 & Estudios sin robot físico: agentes virtuales, chatbots, interfaces, conjuntos de datos o retos de evaluación. \\
CE4 & Estudios secundarios (revisiones), que se emplean solo como antecedentes. \\
CE5 & Registros sin resumen disponible o informes cuyo texto completo no fue accesible. \\
\bottomrule
\end{tabular}
\end{table}

\subsection{Búsqueda sistemática y base de datos de resultados}
La búsqueda se ejecutó el 3 de octubre de 2026 en Semantic Scholar, mediante su API de búsqueda masiva sobre título y resumen (Kinney et al., 2023), y en Crossref, mediante su API REST (Hendricks et al., 2020). Los registros se exportaron a la hoja \textit{Registros} del archivo \textit{Base\_de\_datos.Seminario-Ordaya.xlsx}. Cada registro incluye el identificador, la base de origen, el título, los autores, el año, la fuente, el tipo de documento, el DOI y el resumen, así como la etapa PRISMA, la decisión y el criterio aplicado. La hoja \textit{Ecuaciones} documenta la sintaxis, los filtros, la fecha y el número de registros de cada base, de modo que la búsqueda es reproducible.

\subsection{Proceso de selección}
La selección siguió los pasos de PRISMA 2020 (Page et al., 2021) y se documenta en la Fig.~\ref{fig:prisma}:

\begin{enumerate}
\item \textbf{Identificación.} Se recuperaron 942 registros (873 de Semantic Scholar y 69 de Crossref).
\item \textbf{Eliminación de duplicados.} Se eliminaron 64 duplicados, identificados por coincidencia de DOI o de título normalizado, y quedaron 878 registros únicos.
\item \textbf{Filtros previos al cribado.} Se excluyeron 242 registros por tipo de documento o por tratarse de preprints (CE1) y 44 sin resumen disponible (CE5).
\item \textbf{Cribado por título y resumen.} Se revisaron 592 registros y se excluyeron 497: 82 no abordaban robótica social ni HRI (CE2), 2 no integraban IA generativa (CI2), 231 no describían interacción multimodal (CI3) y 182 no reportaban evaluación empírica ni resultados relacionados con los desenlaces de interés (CI4). El cribado se apoyó en reglas de palabras clave equivalentes a cada criterio, y sus resultados fueron revisados por los autores.
\item \textbf{Recuperación de textos completos.} Se buscó el texto completo de los 95 informes que superaron el cribado. Se usaron las versiones de acceso abierto del editor, Europe PMC, arXiv y repositorios institucionales. Se recuperaron 37 informes. Los 58 restantes no fueron accesibles (28 de IEEE Xplore, 13 de la ACM Digital Library, 8 de MDPI y 9 de otros editores) y se registraron como no recuperados (CE5).
\item \textbf{Elegibilidad a texto completo.} Se leyeron los 37 informes recuperados y se excluyeron 15: 8 sin robot físico (CE3), por ejemplo un audito de sesgos sobre prompts, un dispositivo háptico o conjuntos de datos; 5 con robots sin función social (CE2); y 2 revisiones (CE4). La lectura completa cambió decisiones del cribado. Por ejemplo, Walker et al. (2025) se había descartado por el resumen, pero el texto completo muestra una evaluación con usuarios sobre el robot Pepper.
\item \textbf{Inclusión.} Se incluyeron 22 estudios primarios, publicados en 2023 (1), 2024 (2), 2025 (10) y 2026 (9). La Tabla~\ref{tab:incluidos} resume sus características.
\end{enumerate}

\begin{figure}[htbp]
\centering
\begin{tikzpicture}[
  font=\scriptsize,
  box/.style={draw, rectangle, align=center, text width=4.6cm, minimum height=0.9cm, inner sep=3pt},
  side/.style={draw, rectangle, align=left, text width=5.3cm, minimum height=0.9cm, inner sep=3pt},
  phase/.style={draw, fill=blue!15, rectangle, rotate=90, align=center, minimum width=1.5cm, minimum height=0.5cm, font=\scriptsize\bfseries},
  arr/.style={-{Stealth[length=2mm]}}, node distance=0.55cm]
\node[box] (id) {Registros identificados (n = 942):\\ Semantic Scholar (n = 873)\\ Crossref (n = 69)};
\node[side, right=0.6cm of id] (dup) {Registros eliminados antes del cribado:\\ Duplicados (n = 64)\\ Preprints / tipo no elegible, CE1 (n = 242)\\ Sin resumen, CE5 (n = 44)};
\node[box, below=1.0cm of id] (scr) {Registros cribados por título y resumen\\ (n = 592)};
\node[side, right=0.6cm of scr] (exs) {Registros excluidos (n = 497):\\ CE2 sin robótica social (n = 82)\\ CI2 sin IA generativa (n = 2)\\ CI3 sin interacción multimodal (n = 231)\\ CI4 sin evaluación empírica (n = 182)};
\node[box, below=1.0cm of scr] (rt) {Informes buscados para recuperación\\ (n = 95)};
\node[side, right=0.6cm of rt] (nr) {Informes no recuperados, CE5 (n = 58):\\ IEEE Xplore (n = 28), ACM (n = 13),\\ MDPI (n = 8), otros (n = 9)};
\node[box, below=1.0cm of rt] (el) {Informes evaluados a texto completo\\ (n = 37)};
\node[side, right=0.6cm of el] (exe) {Informes excluidos (n = 15):\\ CE3 sin robot físico (n = 8)\\ CE2 robot sin función social (n = 5)\\ CE4 revisiones (n = 2)};
\node[box, below=1.0cm of el, fill=green!10] (inc) {Estudios incluidos en la RSL\\ (n = 22)};
\draw[arr] (id) -- (scr); \draw[arr] (scr) -- (rt); \draw[arr] (rt) -- (el); \draw[arr] (el) -- (inc);
\draw[arr] (id) -- (dup); \draw[arr] (scr) -- (exs); \draw[arr] (rt) -- (nr); \draw[arr] (el) -- (exe);
\node[phase] at ($(id.west)+(-0.55,0)$) {Identificación};
\node[phase] at ($(scr.west)+(-0.55,0)$) {Cribado};
\node[phase, minimum width=3.4cm] at ($(rt.west)!0.5!(el.west)+(-0.55,0)$) {Elegibilidad};
\node[phase] at ($(inc.west)+(-0.55,0)$) {Inclusión};
\end{tikzpicture}
\caption{Diagrama de flujo PRISMA 2020 del proceso de selección (adaptado de Page et al., 2021).}
\label{fig:prisma}
\end{figure}

\subsection{Extracción de datos}
Los datos de los 22 estudios incluidos se extrajeron de la lectura a texto completo y se registraron en el formulario \textit{Formulario.Seminario-Ordaya.xlsx}, entregado como documento aparte. Cada fila corresponde a un estudio, y las columnas siguen las tres dimensiones del objetivo:
\begin{itemize}
\item datos bibliográficos y de diseño: referencia, año, tipo de estudio, plataforma robótica, población y número de participantes;
\item dimensión técnica: modelo generativo, despliegue (nube, local o híbrido), modalidades de entrada y de salida, arquitectura de fusión, traducción de lenguaje a acción, latencia reportada y su mitigación, sincronización verbal--no verbal y robustez del reconocimiento de voz;
\item dimensiones cognitiva y ética: limitaciones de razonamiento, riesgos éticos identificados y salvaguardas propuestas;
\item cierre: métricas, resultados principales, limitaciones declaradas y dimensión abordada.
\end{itemize}

La Tabla~\ref{tab:incluidos} resume las características principales. Los estudios usan sobre todo Furhat (4), Pepper (4) y ARI (2), y siete se dirigen a adultos mayores. El modelo se despliega en la nube en 10 estudios, en local en 5 y de forma híbrida en 7. Siete estudios reportan tiempos cuantitativos. La latencia de respuesta conversacional va de 0,6~s con predicción de turnos (Skantze e Irfan, 2025) a 4,7~s tras optimizar el pipeline de voz, desde 9,8~s (Satake et al., 2026). La dimensión ética aparece en 11 estudios y la cognitiva en 5. Ningún estudio aborda las tres dimensiones a la vez, lo que confirma el vacío que motiva esta revisión.

{\scriptsize
\setlength{\tabcolsep}{3pt}
\begin{longtable}{@{}r L{1.6cm} L{1.7cm} L{0.9cm} L{2.0cm} L{2.9cm} L{1.4cm}@{}}
\caption{Características de los estudios incluidos (n = 22).}\label{tab:incluidos}\\
\toprule
\textbf{N.º} & \textbf{Estudio} & \textbf{Plataforma} & \textbf{n} & \textbf{Modelo} & \textbf{Hallazgo principal} & \textbf{Dimen\-sión} \\
\midrule
\endfirsthead
\toprule
\textbf{N.º} & \textbf{Estudio} & \textbf{Plataforma} & \textbf{n} & \textbf{Modelo} & \textbf{Hallazgo principal} & \textbf{Dimen\-sión} \\
\midrule
\endhead
\bottomrule
\endfoot
1 & Abbo et al., 2026 & Expressive Furhat: Generating Real-Time Facial Expressions for Human-Robot Dialogue with LLMs & Companion Proceedings of the 21st ACM/IEEE Internati... \\
2 & Abbo et al., 2025b & Multi-party open-ended conversation with a social robot & Frontiers in Robotics and AI \\
3 & Ashok et al., 2025 & ``Thanks for the Practice!'': LLM-Powered Social Robot as Tandem Language Partner at University & 2025 20th ACM/IEEE International Conference on Human... \\
4 & Banna et al., 2025 & Beyond Words: Integrating Personality Traits and Context-Driven Gestures in Human-Robot Interactions & Adaptive Agents and Multi-Agent Systems \\
5 & Bezerra et al., 2026 & Multimodal Navigation Assistance for Older Adults: Combining Generative AI and Augmented Reality in a Robotic Walker & 2026 IEEE 5th International Conference on Intelligen... \\
6 & Blavette et al., 2025 & Acceptability and Usability of a Socially Assistive Robot Integrated With a Large Language Model for Enhanced Human-Robot Interaction in a Geriatric Care Institution: Mixed Methods Evaluation & JMIR Human Factors \\
7 & Blavette et al., 2025 & Integrating a Large Language Model Into a Socially Assistive Robot in a Hospital Geriatric Unit: Two-Wave Comparative Study on Performance, Engagement, and User Perceptions & JMIR Human Factors \\
8 & Chen y Bazzocchi, 2026 & A Dash of Pepper, A Dose of Care: A Context-Aware, Emotionally Adaptive, Intelligent Social Robot for Long-Term Medication Adherence & 2026 11th IEEE RAS/EMBS International Conference on ... \\
9 & Chen et al., 2026 & A Multimodal Emotional Interaction Framework Driven by Large Language Models & 2026 IEEE/ASME International Conference on Advanced ... \\
10 & Chen et al., 2025 & Improving Social Robot's Emotional Intelligence with Physiological Feedback & 2025 IEEE International Conference on Robotics and B... \\
11 & Galatolo y Winkle, 2025 & Simultaneous text and gesture generation for social robots with small language models & Frontiers in Robotics and AI \\
12 & Halima et al., 2026 & Edge AI-Based Object Detection via Voice Recognition with an LLM-Based Emotional Assistant for Elderly Care Robots & International Journal of Information Technology and ... \\
13 & Hari et al., 2025 & Enhancing Student Motivation and Engagement Through the Use of a Slovenian-Speaking Social Robot AlphaMini & Education Sciences \\
14 & Hielscher et al., 2026 & Context-Aware Generation and Modulation of Expressive Motion Behavior using Multimodal Foundation Models & Proceedings of the 21st ACM/IEEE International Confe... \\
15 & Hundt et al., 2024 & LLM-Driven Robots Risk Enacting Discrimination, Violence, and Unlawful Actions & International Journal of Social Robotics \\
16 & Kang et al., 2026 & Fluid-Xpress: Emotion-Aware Dual-Loop Framework for Empathic Facial Reaction in HRI & Companion Proceedings of the 21st ACM/IEEE Internati... \\
17 & Kim et al., 2024 & Understanding Large-Language Model (LLM)-powered Human-Robot Interaction & 2024 19th ACM/IEEE International Conference on Human... \\
18 & Lai et al., 2025 & FAM-HRI: Foundation-Model Assisted Multimodal Human--Robot Interaction Combining Gaze and Speech & IEEE Transactions on Automation Science and Engineer... \\
19 & Lei et al., 2026 & Emotionally intelligent social robot for dementia care: empathy-based conversational intervention model using multimodal sentiment analysis & Future Technology \\
20 & Leusmann et al., 2025 & Investigating LLM-Driven Curiosity in Human-Robot Interaction & Proceedings of the 2025 CHI Conference on Human Fact... \\
21 & Li et al., 2026 & Facilitating Co-regulation: An LLM-Powered Multimodal Social Robot for Parent-Child Dyads & Companion Proceedings of the 21st ACM/IEEE Internati... \\
22 & Liao et al., 2026 & Multimodal Depression Detection Through Conversational Interactions with an Emotion-Aware Social Robot: Pilot Study & JMIR Formative Research \\
23 & Lim et al., 2025 & Design and Implementation of a Companion Robot with LLM-Based Hierarchical Emotion Motion Generation & Applied Sciences \\
24 & Liu et al., 2026 & Multimodal-Language-Model--Driven Interaction and Companionship for Service Robots in Elderly-Care Facilities & 2026 IEEE/ASME International Conference on Advanced ... \\
25 & Liu et al., 2025 & Multimodal Emotion Fusion Mechanism and Empathetic Responses in Companion Robots & IEEE Transactions on Cognitive and Developmental Sys... \\
26 & Lyu et al., 2026 & Enhancing End-user Engagement in Human--Robot Interaction by Performing LLM-driven Expressive Behaviors & ACM Transactions on Human-Robot Interaction \\
27 & Mishra et al., 2023 & Real-time emotion generation in human-robot dialogue using large language models & Frontiers in Robotics and AI \\
28 & Nakagawa y Nihei, 2025 & Bridging Divides through Empathetic Robot Dialogue in Remote Interpersonal Conflict & 2025 34th IEEE International Conference on Robot and... \\
29 & Ok et al., 2026 & Practical Human--Robot Interaction (HRI) through Large Language Model (LLM)-based Voice-to-Action Systems & International Journal of Precision Engineering and M... \\
30 & Pinto y Belpaeme, 2025 & Speech Recognition and LLM Performance in Elderly Care Home Conversations & 2025 34th IEEE International Conference on Robot and... \\
31 & Pinto-Bernal et al., 2025 & Designing Social Robots with LLMs for Engaging Human Interaction & Applied Sciences \\
32 & Pozzi et al., 2026 & Personalization and Generative Dialogue in Social Robotics for Eldercare: A User Study & Applied Sciences \\
33 & Ranasinghe et al., 2025 & Enhancing Receptionist Roles with AI: A Standalone Edge Device based Humanoid Robot Receptionist & 2025 10th International Conference on Control and Ro... \\
34 & Ren y Belpaeme, 2025 & Touched by ChatGPT: Using an LLM to Drive Affective Tactile Interaction & 2025 20th ACM/IEEE International Conference on Human... \\
35 & Satake et al., 2026 & Focus Group-Led Refinement of an LLM-Enabled Companion Robot for Older People & International Journal of Social Robotics \\
36 & Sevilla-Salcedo et al., 2023 & Using Large Language Models to Shape Social Robots' Speech & Int. J. Interact. Multim. Artif. Intell. \\
37 & Skantze y Irfan, 2025 & Applying General Turn-Taking Models to Conversational Human-Robot Interaction & 2025 20th ACM/IEEE International Conference on Human... \\
38 & Souza, 2026 & Integration of Local Generative Language Models in Humanoid Robots: An Approach with the NAO Robot and DeepSeek-LLM & 2026 Brazilian Conference on Robotics (CROS) \\
39 & Steinhaeusser et al., 2026 & Development of the Fully Automatic Robotic Storyteller---Comparing Manual, Semi-Automatic, and LLM-Based Emotional Story Annotation for Multimodal Robotic Storytelling & IEEE Transactions on Affective Computing \\
40 & Syjan et al., 2026 & AI-Driven Humanoid Animatronic Face for Real-Time Multimodal Human-Robot Interaction & 2026 IEEE International Conference on Emerging Syner... \\
41 & Tan et al., 2026 & DR-TSLM: Deictic Reference Resolution via Temporal-Spatial Large Language Model for Human-Robot Interaction & 2026 IEEE International Conference on Cybernetics an... \\
42 & Thabet et al., 2025 & A Unified Framework Low Cost AI-Powered Humanoid Head for Expressive Multimodal Human-Robot Interaction & 2025 International Conference on Future Telecommunic... \\
43 & Viet et al., 2025 & Development of a Humanoid Robot Prototype for Multimodal Human-Robot Interaction & 2025 RIVF International Conference on Computing and ... \\
44 & Wang et al., 2024 & Ain't Misbehavin' - Using LLMs to Generate Expressive Robot Behavior in Conversations with the Tabletop Robot Haru & Companion of the 2024 ACM/IEEE International Confere... \\
45 & Yang et al., 2026 & A Multimodal Adaptive Framework for Social Interaction with the MiRo-E Robot & Sensors (Basel, Switzerland) \\
46 & Zhong et al., 2025 & Integrating LLM into a Socially Assistive Robot for Social Dialogue: An Exploratory Study in a Nursing Home* & 2025 34th IEEE International Conference on Robot and... \\

\end{longtable}
}

%=====================================================================
\section*{Referencias}
%=====================================================================
\begingroup
\renewcommand{\section}[2]{}
\begin{thebibliography}{99}
\small

\bibitem{abbo2025} Abbo, G. A., Desideri, G., Belpaeme, T., y Spitale, M. (2025). ``Can you be my mum?'': Manipulating social robots in the large language models era. En \textit{Proceedings of the 2025 ACM/IEEE International Conference on Human-Robot Interaction (HRI)} (pp. 1181--1185). IEEE.

\bibitem{abercrombie2023} Abercrombie, G., Cercas Curry, A., Dinkar, T., Rieser, V., y Talat, Z. (2023). Mirages: On anthropomorphism in dialogue systems. En \textit{Proceedings of the 2023 Conference on Empirical Methods in Natural Language Processing} (pp. 4776--4790). \url{https://doi.org/10.18653/v1/2023.emnlp-main.290}

\bibitem{barakat2024} Barakat, H., Turk, O., y Demiroglu, C. (2024). Deep learning-based expressive speech synthesis: A systematic review of approaches, challenges, and resources. \textit{EURASIP Journal on Audio, Speech, and Music Processing, 2024}(1), 1--34. \url{https://doi.org/10.1186/s13636-024-00329-7}

\bibitem{hendricks2020} Hendricks, G., Tkaczyk, D., Lin, J., y Feeney, P. (2020). Crossref: The sustainable source of community-owned scholarly metadata. \textit{Quantitative Science Studies, 1}(1), 414--427. \url{https://doi.org/10.1162/qss_a_00022}

\bibitem{kim2024} Kim, C. Y., Lee, C. P., y Mutlu, B. (2024). Understanding large-language model (LLM)-powered human-robot interaction. En \textit{Proceedings of the 2024 ACM/IEEE International Conference on Human-Robot Interaction} (pp. 371--380). \url{https://doi.org/10.1145/3610977.3634966}

\bibitem{kinney2023} Kinney, R., Anastasiades, C., Authur, R., Beltagy, I., Bragg, J., Buraczynski, A., \ldots{} Weld, D. S. (2023). \textit{The Semantic Scholar Open Data Platform} (arXiv:2301.10140). arXiv. \url{https://doi.org/10.48550/arXiv.2301.10140}

\bibitem{kitchenham2007} Kitchenham, B., y Charters, S. (2007). \textit{Guidelines for performing systematic literature reviews in software engineering} (EBSE Technical Report EBSE-2007-01). Keele University y Durham University.

\bibitem{lai2025} Lai, Y., Yuan, S., Nassar, Y., Fan, M., Weber, T., y Rätsch, M. (2025). NVP-HRI: Zero shot natural voice and posture-based human--robot interaction via large language model. \textit{Expert Systems with Applications, 268}, 126360. \url{https://doi.org/10.1016/j.eswa.2025.126360}

\bibitem{maure2025} Maure, R., y Bruno, B. (2025). Autonomy in socially assistive robotics: A systematic review. \textit{Frontiers in Robotics and AI, 12}, 1586473. \url{https://doi.org/10.3389/frobt.2025.1586473}

\bibitem{mauliana2025} Mauliana, M., Ashok, A., Czernochowski, D., y Berns, K. (2025). Exploring LLM-powered multi-session human-robot interactions with university students. \textit{Frontiers in Robotics and AI, 12}, 1585589. \url{https://doi.org/10.3389/frobt.2025.1585589}

\bibitem{page2021} Page, M. J., McKenzie, J. E., Bossuyt, P. M., Boutron, I., Hoffmann, T. C., Mulrow, C. D., \ldots{} Moher, D. (2021). The PRISMA 2020 statement: An updated guideline for reporting systematic reviews. \textit{BMJ, 372}, n71. \url{https://doi.org/10.1136/bmj.n71}

\bibitem{santos2007} Santos, C. M. C., Pimenta, C. A. M., y Nobre, M. R. C. (2007). The PICO strategy for the research question construction and evidence search. \textit{Revista Latino-Americana de Enfermagem, 15}(3), 508--511. \url{https://doi.org/10.1590/S0104-11692007000300023}

\bibitem{verhelst2025} Verhelst, E., y Belpaeme, T. (2025). Large language models cover for speech recognition mistakes: Evaluating conversational AI for second language learners. En \textit{Proceedings of the 2025 ACM/IEEE International Conference on Human-Robot Interaction (HRI)} (pp. 1705--1709). IEEE.

\bibitem{verma2024} Verma, M., Bhambri, S., y Kambhampati, S. (2024). Theory of mind abilities of large language models in human-robot interaction: An illusion? En \textit{Companion of the 2024 ACM/IEEE International Conference on Human-Robot Interaction} (pp. 37--45). \url{https://doi.org/10.1145/3610978.3640767}

\bibitem{zhang2023} Zhang, B., y Soh, H. (2023). Large language models as zero-shot human models for human-robot interaction. En \textit{2023 IEEE/RSJ International Conference on Intelligent Robots and Systems (IROS)} (pp. 7961--7968). \url{https://doi.org/10.1109/IROS55552.2023.10341488}

\end{thebibliography}
\endgroup

\subsection*{Estudios incluidos en la RSL}
\begingroup
\renewcommand{\section}[2]{}
\begin{thebibliography}{99}
\small
\bibitem{R0070} Abbo, G. A., Janssens, R., Vreken, S. V. d., y Belpaeme, T. (2026). Expressive Furhat: Generating Real-Time Facial Expressions for Human-Robot Dialogue with LLMs. \textit{Companion Proceedings of the 21st ACM/IEEE International Conference on Human-Robot Interaction}. \url{https://doi.org/10.1145/3776734.3794436}

\bibitem{R0617} Abbo, G. A., Pinto-Bernal, M. J., Catrycke, M., y Belpaeme, T. (2025b). Multi-party open-ended conversation with a social robot. \textit{Frontiers in Robotics and AI}. \url{https://doi.org/10.3389/frobt.2026.1766383}

\bibitem{R0292} Ashok, A., Bruno, B., Helf, T., y Berns, K. (2025). ``Thanks for the Practice!'': LLM-Powered Social Robot as Tandem Language Partner at University. \textit{2025 20th ACM/IEEE International Conference on Human-Robot Interaction (HRI)}. \url{https://doi.org/10.1109/hri61500.2025.10973837}

\bibitem{R0588} Banna, T. T., Rahman, S., y Tareq, M. (2025). Beyond Words: Integrating Personality Traits and Context-Driven Gestures in Human-Robot Interactions. \textit{Adaptive Agents and Multi-Agent Systems}. \url{https://doi.org/10.5555/3709347.3743537}

\bibitem{R0488} Bezerra, M., Loureiro, M., Machado, F., Rocon, E., y Frizera, A. (2026). Multimodal Navigation Assistance for Older Adults: Combining Generative AI and Augmented Reality in a Robotic Walker. \textit{2026 IEEE 5th International Conference on Intelligent Reality (ICIR)}. \url{https://doi.org/10.1109/icir70631.2026.11628389}

\bibitem{R0623} Blavette, L., Dacunha, S., Alameda-Pineda, X., García, D. H., Gannot, S., Gras, F., Gunson, N., Lemaignan, S., Polic, M., Tandeitnik, P., Tonini, F., Rigaud, A., y Pino, M. (2025). Acceptability and Usability of a Socially Assistive Robot Integrated With a Large Language Model for Enhanced Human-Robot Interaction in a Geriatric Care Institution: Mixed Methods Evaluation. \textit{JMIR Human Factors}. \url{https://doi.org/10.2196/76496}

\bibitem{R0671} Blavette, L., Dacunha, S., Alameda-Pineda, X., Cattoni, J., Rigaud, A., y Pino, M. (2025). Integrating a Large Language Model Into a Socially Assistive Robot in a Hospital Geriatric Unit: Two-Wave Comparative Study on Performance, Engagement, and User Perceptions. \textit{JMIR Human Factors}. \url{https://doi.org/10.2196/81936}

\bibitem{R0072} Chen, A., y Bazzocchi, M. C. F. (2026). A Dash of Pepper, A Dose of Care: A Context-Aware, Emotionally Adaptive, Intelligent Social Robot for Long-Term Medication Adherence. \textit{2026 11th IEEE RAS/EMBS International Conference on Biomedical Robotics and Biomechatronics (BioRob)}. \url{https://doi.org/10.1109/biorob66782.2026.11681492}

\bibitem{R0473} Chen, Y., Qiu, C., Wang, N., Cao, J., y Xiong, Z. (2026). A Multimodal Emotional Interaction Framework Driven by Large Language Models. \textit{2026 IEEE/ASME International Conference on Advanced Intelligent Mechatronics (AIM)}. \url{https://doi.org/10.1109/aim65483.2026.11658264}

\bibitem{R0479} Chen, Y., Hu, M., Tian, L., Nichols, E., Gomez, R., y Li, G. (2025). Improving Social Robot's Emotional Intelligence with Physiological Feedback. \textit{2025 IEEE International Conference on Robotics and Biomimetics (ROBIO)}. \url{https://doi.org/10.1109/robio66223.2025.11377720}

\bibitem{R0445} Galatolo, A., y Winkle, K. (2025). Simultaneous text and gesture generation for social robots with small language models. \textit{Frontiers in Robotics and AI}. \url{https://doi.org/10.3389/frobt.2025.1581024}

\bibitem{R0477} Halima, S. B., Abdallah, F. B., y Haggège, J. (2026). Edge AI-Based Object Detection via Voice Recognition with an LLM-Based Emotional Assistant for Elderly Care Robots. \textit{International Journal of Information Technology and Computer Science}. \url{https://doi.org/10.5815/ijitcs.2026.04.04}

\bibitem{R0516} Hari, D., Skrbinjek, V., y Flogie, A. (2025). Enhancing Student Motivation and Engagement Through the Use of a Slovenian-Speaking Social Robot AlphaMini. \textit{Education Sciences}. \url{https://doi.org/10.3390/educsci15091222}

\bibitem{R0697} Hielscher, T., Scaparro, F., y Arras, K. O. (2026). Context-Aware Generation and Modulation of Expressive Motion Behavior using Multimodal Foundation Models. \textit{Proceedings of the 21st ACM/IEEE International Conference on Human-Robot Interaction}. \url{https://doi.org/10.1145/3757279.3785635}

\bibitem{R0423} Hundt, A., Azeem, R., Mansouri, M., y Brandão, M. (2024). LLM-Driven Robots Risk Enacting Discrimination, Violence, and Unlawful Actions. \textit{International Journal of Social Robotics}. \url{https://doi.org/10.1007/s12369-025-01301-x}

\bibitem{R0416} Kang, C., Alitai, M., Wang, Y., Cai, X., Ma, R., Cangelosi, A., y Shangguan, Z. (2026). Fluid-Xpress: Emotion-Aware Dual-Loop Framework for Empathic Facial Reaction in HRI. \textit{Companion Proceedings of the 21st ACM/IEEE International Conference on Human-Robot Interaction}. \url{https://doi.org/10.1145/3776734.3794429}

\bibitem{R0210} Kim, C. Y., Lee, C. P., y Mutlu, B. (2024). Understanding Large-Language Model (LLM)-powered Human-Robot Interaction. \textit{2024 19th ACM/IEEE International Conference on Human-Robot Interaction (HRI)}. \url{https://doi.org/10.1145/3610977.3634966}

\bibitem{R0556} Lai, Y., Yuan, S., Li, P., Zhang, B., Kiefer, B., Deng, T., y Zell, A. (2025). FAM-HRI: Foundation-Model Assisted Multimodal Human--Robot Interaction Combining Gaze and Speech. \textit{IEEE Transactions on Automation Science and Engineering}. \url{https://doi.org/10.1109/tase.2026.3695438}

\bibitem{R0779} Lei, Z., Yin, Y., y Chen, Y. (2026). Emotionally intelligent social robot for dementia care: empathy-based conversational intervention model using multimodal sentiment analysis. \textit{Future Technology}. \url{https://doi.org/10.55670/fpll.futech.5.2.19}

\bibitem{R0754} Leusmann, J., Belardinelli, A., Haliburton, L., Hasler, S., Schmidt, A., Mayer, S., Gienger, M., y Wang, C. (2025). Investigating LLM-Driven Curiosity in Human-Robot Interaction. \textit{Proceedings of the 2025 CHI Conference on Human Factors in Computing Systems}. \url{https://doi.org/10.1145/3706598.3713923}

\bibitem{R0286} Li, J., Schijve, F., Hu, J., y Barakova, E. (2026). Facilitating Co-regulation: An LLM-Powered Multimodal Social Robot for Parent-Child Dyads. \textit{Companion Proceedings of the 21st ACM/IEEE International Conference on Human-Robot Interaction}. \url{https://doi.org/10.1145/3776734.3794485}

\bibitem{R0538} Liao, P., Su, Y., Qian, X., Chang, Y., Lee, Y., y Fu, L. (2026). Multimodal Depression Detection Through Conversational Interactions with an Emotion-Aware Social Robot: Pilot Study. \textit{JMIR Formative Research}. \url{https://doi.org/10.2196/84110}

\bibitem{R0852} Lim, Y., Cho, J., Lee, D., Choi, D., y Lee, D. (2025). Design and Implementation of a Companion Robot with LLM-Based Hierarchical Emotion Motion Generation. \textit{Applied Sciences}. \url{https://doi.org/10.3390/app152312759}

\bibitem{R0020} Liu, C., Vu, C., y Liu, Y. (2026). Multimodal-Language-Model--Driven Interaction and Companionship for Service Robots in Elderly-Care Facilities. \textit{2026 IEEE/ASME International Conference on Advanced Intelligent Mechatronics (AIM)}. \url{https://doi.org/10.1109/aim65483.2026.11658058}

\bibitem{R0460} Liu, X., Lv, Q., Li, J., Song, S., y Cangelosi, A. (2025). Multimodal Emotion Fusion Mechanism and Empathetic Responses in Companion Robots. \textit{IEEE Transactions on Cognitive and Developmental Systems}. \url{https://doi.org/10.1109/tcds.2024.3442203}

\bibitem{R0333} Lyu, S., Wang, F., Lin, W., Guo, G., y Navarro-Alarcón, D. (2026). Enhancing End-user Engagement in Human--Robot Interaction by Performing LLM-driven Expressive Behaviors. \textit{ACM Transactions on Human-Robot Interaction}. \url{https://doi.org/10.1145/3813107}

\bibitem{R0817} Mishra, C., Verdonschot, R. G., Hagoort, P., y Skantze, G. (2023). Real-time emotion generation in human-robot dialogue using large language models. \textit{Frontiers in Robotics and AI}. \url{https://doi.org/10.3389/frobt.2023.1271610}

\bibitem{R0264} Nakagawa, S., y Nihei, M. (2025). Bridging Divides through Empathetic Robot Dialogue in Remote Interpersonal Conflict. \textit{2025 34th IEEE International Conference on Robot and Human Interactive Communication (RO-MAN)}. \url{https://doi.org/10.1109/ro-man63969.2025.11217560}

\bibitem{R0185} Ok, K., Heo, G., Lee, C., y Ahn, S. (2026). Practical Human--Robot Interaction (HRI) through Large Language Model (LLM)-based Voice-to-Action Systems. \textit{International Journal of Precision Engineering and Manufacturing}. \url{https://doi.org/10.1007/s12541-026-01573-x}

\bibitem{R0492} Pinto, M. J., y Belpaeme, T. (2025). Speech Recognition and LLM Performance in Elderly Care Home Conversations. \textit{2025 34th IEEE International Conference on Robot and Human Interactive Communication (RO-MAN)}. \url{https://doi.org/10.1109/ro-man63969.2025.11217575}

\bibitem{R0212} Pinto-Bernal, M. J., Biondina, M., y Belpaeme, T. (2025). Designing Social Robots with LLMs for Engaging Human Interaction. \textit{Applied Sciences}. \url{https://doi.org/10.3390/app15116377}

\bibitem{R0400} Pozzi, L., Nasato, M., Toscani, N., Braghin, F., y Gandolla, M. (2026). Personalization and Generative Dialogue in Social Robotics for Eldercare: A User Study. \textit{Applied Sciences}. \url{https://doi.org/10.3390/app16073369}

\bibitem{R0615} Ranasinghe, C., Munasinghe, V., Vikasitha, S., Abeywickrama, C., Silva, B., y Jayasekara, P. (2025). Enhancing Receptionist Roles with AI: A Standalone Edge Device based Humanoid Robot Receptionist. \textit{2025 10th International Conference on Control and Robotics Engineering (ICCRE)}. \url{https://doi.org/10.1109/iccre65455.2025.11093535}

\bibitem{R0431} Ren, Q., y Belpaeme, T. (2025). Touched by ChatGPT: Using an LLM to Drive Affective Tactile Interaction. \textit{2025 20th ACM/IEEE International Conference on Human-Robot Interaction (HRI)}. \url{https://doi.org/10.1109/hri61500.2025.10973852}

\bibitem{R0076} Satake, Y., Nygren, M., Yu, C., Tsuboi, A., Endo, I., Umemura, K., Katsuki, K., Ishimaru, D., Rapaport, P., Bianchi-Berthouze, N., Ikeda, M., y Howard, R. (2026). Focus Group-Led Refinement of an LLM-Enabled Companion Robot for Older People. \textit{International Journal of Social Robotics}. \url{https://doi.org/10.1007/s12369-026-01407-w}

\bibitem{R0862} Sevilla-Salcedo, J., Fernández-Rodicio, E., Martín-Galván, L., González, Á. C., Castillo, J., y Salichs, M. (2023). Using Large Language Models to Shape Social Robots' Speech. \textit{Int. J. Interact. Multim. Artif. Intell.}. \url{https://doi.org/10.9781/ijimai.2023.07.008}

\bibitem{R0051} Skantze, G., y Irfan, B. (2025). Applying General Turn-Taking Models to Conversational Human-Robot Interaction. \textit{2025 20th ACM/IEEE International Conference on Human-Robot Interaction (HRI)}. \url{https://doi.org/10.1109/hri61500.2025.10973958}

\bibitem{R0528} Souza, V. (2026). Integration of Local Generative Language Models in Humanoid Robots: An Approach with the NAO Robot and DeepSeek-LLM. \textit{2026 Brazilian Conference on Robotics (CROS)}. \url{https://doi.org/10.1109/cros69211.2026.11565712}

\bibitem{R0575} Steinhaeusser, S. C., Zehe, A., Maier, S., Hotho, A., y Lugrin, B. (2026). Development of the Fully Automatic Robotic Storyteller---Comparing Manual, Semi-Automatic, and LLM-Based Emotional Story Annotation for Multimodal Robotic Storytelling. \textit{IEEE Transactions on Affective Computing}. \url{https://doi.org/10.1109/taffc.2026.3691816}

\bibitem{R0226} Syjan, S., Jaradath, H. T., y Shetty, V. C. (2026). AI-Driven Humanoid Animatronic Face for Real-Time Multimodal Human-Robot Interaction. \textit{2026 IEEE International Conference on Emerging Synergy Science and Technology (ICESST)}. \url{https://doi.org/10.1109/icesst69086.2026.11582887}

\bibitem{R0688} Tan, R., Lou, S., Yu, L., Tian, X., y Lv, C. (2026). DR-TSLM: Deictic Reference Resolution via Temporal-Spatial Large Language Model for Human-Robot Interaction. \textit{2026 IEEE International Conference on Cybernetics and Intelligent Systems (CIS) and IEEE International Conference on Robotics, Automation and Mechatronics (RAM)}. \url{https://doi.org/10.1109/cis-ram71614.2026.11668533}

\bibitem{R0591} Thabet, M. M., Gad, E., Abdelmonem, H. K., y Ibrahim, R. M. (2025). A Unified Framework Low Cost AI-Powered Humanoid Head for Expressive Multimodal Human-Robot Interaction. \textit{2025 International Conference on Future Telecommunications and Artificial Intelligence (IC-FTAI)}. \url{https://doi.org/10.1109/ic-ftai67960.2025.11384544}

\bibitem{R0201} Viet, T. T., Canh, T., Gia, H. U., Van, P. D., Duc, S. T., Do, N. M., y HoangVan, X. (2025). Development of a Humanoid Robot Prototype for Multimodal Human-Robot Interaction. \textit{2025 RIVF International Conference on Computing and Communication Technologies (RIVF)}. \url{https://doi.org/10.1109/rivf68649.2025.11364526}

\bibitem{R0236} Wang, Z., Reisert, P., Nichols, E., y Gomez, R. (2024). Ain't Misbehavin' - Using LLMs to Generate Expressive Robot Behavior in Conversations with the Tabletop Robot Haru. \textit{Companion of the 2024 ACM/IEEE International Conference on Human-Robot Interaction}. \url{https://doi.org/10.1145/3610978.3640562}

\bibitem{R0293} Yang, Y., Yap, P. S., Wijeakumar, S., Magassouba, A., y Deshpande, N. (2026). A Multimodal Adaptive Framework for Social Interaction with the MiRo-E Robot. \textit{Sensors (Basel, Switzerland)}. \url{https://doi.org/10.3390/s26041209}

\bibitem{R0405} Zhong, V. J., Studerus, E., y Vonschallen, S. (2025). Integrating LLM into a Socially Assistive Robot for Social Dialogue: An Exploratory Study in a Nursing Home*. \textit{2025 34th IEEE International Conference on Robot and Human Interactive Communication (RO-MAN)}. \url{https://doi.org/10.1109/ro-man63969.2025.11217629}

\end{thebibliography}
\endgroup

\end{document}
